% ============================================================
%  DOCUMENTO DE ESTUDIO PARA LA DEFENSA
%  Proyecto: Efficient Fine-Tuning of Small Open LLMs for
%            Spanish Sentiment Analysis (LoRA/QLoRA)
%  Autor del proyecto: Beltrán Sánchez Careaga
%
%  Material de estudio (NO es el paper). Por cada concepto:
%  definición rigurosa primero, intuición/analogía como apoyo,
%  procedimiento, caja de análisis (qué/porqué/alternativas/
%  qué NO significa) y autotest. En los ítems difíciles, el
%  autotest incluye una "derivación en blanco".
%
%  Estado: BLOQUE 1 COMPLETO (14 ítems).
% ============================================================

\documentclass[11pt,a4paper]{article}

\usepackage[utf8]{inputenc}
\usepackage[T1]{fontenc}   % acentos y guiones correctos sin depender de spanish.ldf
% Nota: se evita babel-spanish porque spanish.ldf no está instalado en este entorno
% (solo el locale .ini). El contenido es español; la tipografía no lo requiere.
\usepackage{lmodern}
\usepackage{amsmath}
\usepackage{amssymb}
\usepackage{booktabs}
\usepackage{geometry}
\geometry{margin=2.5cm}
\usepackage{xcolor}
\usepackage{enumitem}
\usepackage{hyperref}
\hypersetup{colorlinks=true, linkcolor=blue!50!black, urlcolor=blue!50!black}

% ---- Comandos de apoyo ----
\newcommand{\nivel}[1]{\medskip\noindent\textbf{\textcolor{blue!45!black}{#1}}\par\nopagebreak\smallskip}
\newcommand{\fuente}[1]{\nopagebreak\par\hfill{\footnotesize\textit{Fuente: #1}}\par}
\newcommand{\externo}[1]{{\footnotesize\textcolor{red!55!black}{[externo: #1]}}}
\newcommand{\notrace}[1]{{\footnotesize\textcolor{orange!75!black}{\textbf{[NO trazable a tabla/sección: #1]}}}}

% Caja "qué / por qué / alternativas / interpretación"
\newenvironment{analisis}{%
  \medskip\par\noindent\rule{\linewidth}{0.4pt}\par\nopagebreak
  \textbf{Análisis de la decisión / resultado}\par\nopagebreak\smallskip
  \begin{description}[leftmargin=2.6cm,style=nextline,font=\bfseries]
}{%
  \end{description}\noindent\rule{\linewidth}{0.4pt}\par\medskip
}

\title{\textbf{Documento de estudio para la defensa}\\[2pt]
\large Fine-Tuning eficiente (LoRA/QLoRA) de LLMs pequeños para análisis de sentimiento en español}
\author{Material de preparación --- Bloque 1: conceptos densos}
\date{}

\begin{document}
\maketitle

% ============================================================
\section*{Cómo usar este documento}
% ============================================================
Cada concepto se explica con la \textbf{definición rigurosa primero} y la \textbf{intuición /
analogía como segundo apoyo}, seguidas del \textbf{procedimiento}. Después, una caja de
\textbf{análisis} (qué se decidió/observó, por qué, qué se descartó y, obligatoriamente, qué
significa y \emph{qué NO significa}) y un \textbf{autotest}. En los ítems difíciles el autotest
termina con una \emph{derivación en blanco}: con solo las dimensiones de entrada, reconstruir la
fórmula y la cuenta de parámetros de memoria.

\medskip\noindent Convenciones de fidelidad a la fuente:
\begin{itemize}[nosep]
  \item Todo número está trazado a \texttt{paper/main.tex} (sección \S\ o Tabla I--V).
  \item El conocimiento externo al paper va en rojo: \externo{ejemplo}.
  \item Cualquier número \emph{no} trazable a una tabla/sección va marcado en naranja:
        \notrace{ejemplo ilustrativo}.
\end{itemize}

% ============================================================
\section{Cómo leer este proyecto desde cero}
% ============================================================

\nivel{El problema, en una frase}
Hay millones de tuits en español con opiniones, y queremos que una máquina diga si cada uno es
\emph{positivo}, \emph{negativo} o \emph{neutral}. El problema no es solo acertar: es acertar
\emph{barato}.

\nivel{Por qué era caro hasta ahora}
Los modelos de lenguaje grandes (LLMs) son como un cerebro con miles de millones de \emph{pesos}
(números que codifican lo que el modelo «sabe»). Adaptar ese cerebro a una tarea concreta, a la
manera clásica, significa reescribir todos esos pesos (\emph{full fine-tuning}): hace falta una
GPU cara y muchos datos etiquetados. Para un estudiante, un autónomo o una pyme, eso es
prohibitivo.

\nivel{La idea que lo cambia: PEFT (LoRA / QLoRA)}
En vez de reescribir el cerebro entero, lo \textbf{congelamos} y le pegamos unas «notas adhesivas»
diminutas que corrigen su comportamiento. Eso es \textbf{LoRA}: se entrenan solo unas matrices
pequeñas, una fracción mínima del modelo. \textbf{QLoRA} va más lejos: además comprime el cerebro
congelado a 4 bits para que ocupe mucha menos memoria de GPU. Resultado: se puede afinar un modelo
decente en una GPU \emph{gratuita} (el objetivo del proyecto es una T4 de 16~GB).

\nivel{El reparto de personajes}
\begin{itemize}[nosep]
  \item \textbf{Modelos a afinar}: Qwen3-1.7B y Qwen3-4B (LLMs generativos abiertos, pequeños).
  \item \textbf{Rivales clásicos} (\emph{baselines}): BETO y XLM-RoBERTa (encoders de español,
        afinados a la antigua), más un suelo léxico (TF-IDF) y una referencia de fábrica
        (RoBERTuito).
  \item \textbf{Datos}: tuits en español de Cardiff~NLP (cortes A y B); e InterTASS~2018, con
        tuits de España, Costa Rica y Perú (corte~C).
  \item \textbf{La regla de medir}: \emph{macro-F1}, una nota que pondera por igual las tres clases.
\end{itemize}

\nivel{Las tres preguntas (los tres «cortes») y sus respuestas en una línea}
\begin{description}[leftmargin=1.6cm,style=nextline,font=\bfseries]
  \item[Cut A] \emph{¿Cuántos ejemplos hacen falta para que afinar gane a solo «pedírselo» (prompting) al modelo?}
        \\ Depende del tamaño: al 1.7B le bastan $<7$ ejemplos; el 4B necesita $\approx 50$, porque de partida ya lo hace mejor sin afinar.
  \item[Cut B] \emph{¿Un modelo pequeño con LoRA gana a los encoders clásicos teniendo en cuenta el coste?}
        \\ Sí. LoRA «domina la frontera»: más calidad (F1 0.707 / 0.698 vs 0.661 de BETO) entrenando una fracción de los parámetros.
  \item[Cut C] \emph{Un adaptador entrenado en una variedad de español, ¿sirve para otras?}
        \\ Casi siempre sí, pero con una asimetría: entrenar en costarricense o peruano y evaluar en español de España cuesta $0.031$--$0.039$ de F1.
\end{description}

\nivel{La moraleja}
El proyecto no vende «un modelo mejor», sino \textbf{evidencia reproducible} sobre \emph{cuándo} y
\emph{a qué coste} conviene cada opción, en hardware al alcance de cualquiera. Por eso muchos
resultados que parecen «negativos» (una penalización aquí, QLoRA más lento allá) son en realidad la
prueba de que el experimento estaba bien diseñado: midió lo que decía medir.

\bigskip\hrule\bigskip

% ============================================================
\section{LoRA: adaptar un modelo gigante moviendo casi nada}\label{sec:lora}
% ============================================================

\nivel{Definición rigurosa}
Sea $W_0 \in \mathbb{R}^{d \times k}$ una matriz de pesos preentrenada y \textbf{congelada}. LoRA
reparametriza su actualización como un producto de \emph{rango bajo}:
\[
\Delta W = B A, \qquad B \in \mathbb{R}^{d \times r}, \quad A \in \mathbb{R}^{r \times k}, \quad r \ll \min(d,k).
\]
La capa, en lugar de $h = W_0 x$, calcula
\[
h = W_0 x + \frac{\alpha}{r}\, B A\, x .
\]
Solo $A$ y $B$ se entrenan; $W_0$ permanece fijo. El factor $\alpha/r$ escala la contribución del
adaptador. En el proyecto: $r=16$, $\alpha=32$ (luego $\alpha/r=2$), \emph{dropout} $0{,}05$, sobre
las proyecciones de atención \texttt{\{q,k,v,o\}\_proj}.
\fuente{\S III-C «LoRA Configuration» de \texttt{paper/main.tex}}

\nivel{La intuición (apoyo)}
Un LLM es como una \textbf{enciclopedia enorme} ya impresa que no queremos reimprimir. LoRA la deja
\textbf{congelada} y le añade una \textbf{hoja de erratas muy fina}: «donde pone esto, entiende esto
otro». Esa hoja (las matrices $A,B$) es pequeñísima porque muchos ajustes útiles «caben» en pocas
dimensiones; el número de dimensiones es el rango $r$ (aquí 16): más rango = corrección más
detallada, pero más cara.

\nivel{El procedimiento}
\begin{enumerate}[nosep]
  \item Eliges qué matrices corregir: aquí \texttt{\{q,k,v,o\}\_proj} (atención).
  \item Congelas cada $W_0$ (no se entrena).
  \item Añades $\Delta W = BA$ y entrenas \emph{solo} $A,B$.
  \item Escalas por $\alpha/r$.
  \item Al inferir, $\Delta W$ se puede \emph{fundir} en $W_0$: en teoría, sin coste extra.
\end{enumerate}

\noindent\textbf{Cuenta de parámetros (el corazón del ahorro).} Entrenar $W_0$ entera cuesta
$d\cdot k$ parámetros; LoRA cuesta $r\,(d+k)$. Con $d=k=2048$ y $r=16$:
\[
\text{full} = 2048^2 \approx 4{,}19\text{M}, \qquad \text{LoRA} = 16\,(2048{+}2048) = 65\,536 \approx 0{,}066\text{M},
\]
es decir $\approx 64\times$ menos \emph{por matriz}.
\notrace{el valor $d=k=2048$ es ilustrativo; el paper no publica la dimensión oculta de Qwen3}.
Sumando todas las matrices adaptadas, el proyecto entrena \textbf{6,42~M} (1.7B) y \textbf{11,80~M}
(4B) parámetros \emph{(Tabla~IV)} sobre modelos de $1{,}72$~B y $4$~B pesos totales.

\medskip\noindent\textbf{Sobre el «$\sim\!10\,000\times$» (debilidad W14).} El factor $10\,000\times$
es del paper original de LoRA \externo{Hu et al., 2021: $r{=}4$ sobre GPT-3 de 175\,B} y así se cita
en \S II-A. En \emph{este} proyecto, comparando parámetros entrenables, lo real es:
\[
\frac{\text{BETO }110\text{M}}{6{,}42\text{M}} \approx 17\times,\quad
\frac{\text{XLM-R }278\text{M}}{6{,}42\text{M}} \approx 43\times,\quad
\frac{\text{full-FT }1{,}72\text{B}}{6{,}42\text{M}} \approx 268\times .
\]
El argumento de fondo (LoRA entrena una fracción minúscula) es verdadero; solo el número
«$10\,000\times$» no aplica a esta comparación. \emph{(Verificado contra Tablas~I y~IV.)}

\begin{analisis}
  \item[Qué] Reparametrización de la actualización como producto de rango bajo $\Delta W=BA$,
        entrenando solo $A,B$ y congelando $W_0$.
  \item[Por qué] Adaptar modelos de miles de millones de pesos entrenando millones; cabe en GPU
        gratuita y, al poder fundir $\Delta W$, en teoría no añade latencia.
  \item[Alternativas descartadas] \emph{Full-FT} (mueve todo, $\approx 270\times$ más, casi doble
        VRAM, y aquí no mejora calidad); \emph{adapters clásicos} \externo{Houlsby et al., 2019}
        (añaden latencia); \emph{DoRA/IA$^3$} (\emph{future work}, no evaluados).
  \item[Qué significa] Con $r=16$ sobre la atención basta para superar a encoders full-fine-tuned
        en esta tarea.
  \item[Qué NO significa] No que LoRA iguale a full-FT en \emph{cualquier} tarea; no
        «$10\,000\times$» en este setup (W14); y el «sin latencia» es cierto \emph{solo si} se funde
        el adaptador (la Tabla~IV mide 54 vs 33~ms porque se evaluó sin fundir).
\end{analisis}

\nivel{Autotest}
\begin{enumerate}[nosep]
  \item ¿Para qué sirve el factor $\alpha/r$ al variar $r$?
  \item ¿Por qué la teoría dice «0 latencia» pero la Tabla~IV mide 54~ms (LoRA) vs 33~ms (full-FT)?
  \item ¿De dónde sale el «$10\,000\times$» y por qué no aplica a tu Cut~B? Da las cifras correctas.
  \item \textbf{Derivación en blanco.} Dados $d=k=2048$ y $r=16$: escribe $\Delta W$ con las
        dimensiones de $A$ y $B$, deriva el número de parámetros entrenables por matriz con LoRA y
        con full-FT, y el cociente. (Solución: $BA$ con $B\in\mathbb{R}^{2048\times16}$,
        $A\in\mathbb{R}^{16\times2048}$; LoRA $=16(2048{+}2048)=65\,536$; full $=2048^2=4\,194\,304$;
        $\approx 64\times$.)
\end{enumerate}

% ============================================================
\section{QLoRA: la misma idea, pero el cerebro comprimido a 4 bits}\label{sec:qlora}
% ============================================================

\nivel{Definición rigurosa}
QLoRA carga el modelo base en \textbf{4 bits} con la \texttt{BitsAndBytesConfig} del proyecto: tipo
de cuantización \texttt{nf4} (\emph{NormalFloat-4}), doble cuantización activada
(\texttt{use\_double\_quant=True}) y \texttt{compute\_dtype = bfloat16}. Tras cargar, se aplica
\texttt{prepare\_model\_for\_kbit\_training} y se activa \emph{gradient checkpointing}; sobre ese
backbone cuantizado se monta la \emph{misma} configuración LoRA ($r=16$, $\alpha=32$, etc.).
\fuente{\S III-C «QLoRA Configuration», \texttt{src/models/lora\_finetune.py}}

\nivel{La intuición (apoyo)}
Si LoRA congela la enciclopedia, QLoRA además la \textbf{imprime en papel finísimo}: cada peso, que
ocupaba 16 bits, se guarda con solo 4 bits (16 niveles posibles en vez de miles). El cerebro pesa
mucho menos en memoria; el coste es que, para «leerlo», hay que reconstruir cada peso al vuelo
(\emph{dequantización}), lo que cuesta tiempo. NF4 es un formato de 4 bits cuyos niveles se reparten
suponiendo que los pesos siguen una campana de Gauss \externo{Dettmers et al., 2023}, de modo que se
pierde menos información que con 4 bits «planos».

\nivel{Resultados clave en el proyecto (Tabla~IV)}
\begin{itemize}[nosep]
  \item \textbf{Memoria}: VRAM pico baja a \textbf{4,93~GB} (1.7B) y \textbf{6,79~GB} (4B), frente a
        $8{,}72$ y $16{,}7$~GB de LoRA: reducción de $\mathbf{1{,}8\times}$ y $\mathbf{2{,}5\times}$.
        El 4B entra así en GPUs de consumo de 8~GB.
  \item \textbf{Calidad}: retiene $0{,}681/0{,}698 = \mathbf{97{,}6\%}$ (1.7B) y
        $0{,}701/0{,}707 = \mathbf{99{,}2\%}$ (4B) de la F1 de LoRA.
  \item \textbf{Coste}: entrena $\approx\mathbf{2{,}4\times}$ más lento (409 vs 169~s; 534 vs 218~s)
        y la latencia de inferencia $\approx\mathbf{2\times}$ (102 vs 54~ms; 132 vs 68~ms), por la
        dequantización al vuelo.
\end{itemize}

\begin{analisis}
  \item[Qué] Cuantización NF4 a 4 bits del backbone congelado + doble cuantización, con LoRA encima.
  \item[Por qué] Reducir la VRAM pico para meter el 4B en hardware de consumo (8~GB), el objetivo
        «free-tier» del proyecto.
  \item[Alternativas descartadas] LoRA sin cuantizar (el 4B no baja de 16~GB); cuantización de 8
        bits (menos ahorro); 4 bits «planos» en vez de NF4 (peor preservación de la información).
  \item[Qué significa] A esta escala, cuantizar a 4 bits cuesta \emph{casi nada} en calidad
        (1--2,5 puntos de F1) y mucho en VRAM. La elección LoRA-vs-QLoRA debe gobernarse por el eje
        de coste (VRAM/latencia), no por la calidad esperada.
  \item[Qué NO significa] No que QLoRA sea «gratis»: paga $2{,}4\times$ en entrenamiento y $2\times$
        en latencia. Para inferencia de baja latencia ($<100$~ms), LoRA sin cuantizar es mejor si la
        VRAM lo permite.
\end{analisis}

\nivel{Autotest}
\begin{enumerate}[nosep]
  \item ¿Por qué QLoRA es más \emph{lento} aunque ocupe menos memoria?
  \item ¿Qué aporta NF4 frente a 4 bits «planos»?
  \item Para el 4B, ¿cuál es la retención de F1 y la reducción de VRAM? ¿Qué implica para elegir
        entre LoRA y QLoRA?
  \item \textbf{Derivación en blanco.} A partir de las cifras de la Tabla~IV, calcula el factor de
        reducción de VRAM de QLoRA vs LoRA para 1.7B y 4B, y el porcentaje de F1 retenido en cada
        caso. (Solución: $8{,}72/4{,}93\approx1{,}8\times$ y $16{,}7/6{,}79\approx2{,}5\times$;
        $97{,}6\%$ y $99{,}2\%$.)
\end{enumerate}

% ============================================================
\section{Macro-F1: la regla con la que se mide todo}\label{sec:macrof1}
% ============================================================

\nivel{Definición rigurosa}
Para una clase $c$, con $TP_c,FP_c,FN_c$ (verdaderos positivos, falsos positivos, falsos negativos):
\[
P_c = \frac{TP_c}{TP_c + FP_c}, \quad
R_c = \frac{TP_c}{TP_c + FN_c}, \quad
F1_c = \frac{2\,P_c\,R_c}{P_c + R_c}\ \text{(media armónica)}.
\]
El \textbf{macro-F1} promedia sin pesos sobre las $C=3$ clases:
\[
\text{macro-}F1 = \frac{1}{C}\sum_{c=1}^{C} F1_c = \frac{F1_{\text{neg}}+F1_{\text{neu}}+F1_{\text{pos}}}{3}.
\]
\fuente{\S III-D y \S IV de \texttt{paper/main.tex}}

\nivel{La intuición (apoyo)}
Es un examen con \textbf{tres temas de igual peso}: tu nota es la \emph{media de las tres notas},
no el total de aciertos. Si fuera el total, podrías aprobar dominando los dos temas fáciles e
ignorando el difícil. Aquí el tema difícil es \textbf{neutral} (detectar que un tuit \emph{no}
expresa emoción). \emph{Accuracy} es «total de aciertos» y puede salir alta ignorando neutral;
macro-F1 obliga a no abandonarlo. La media armónica es alta solo si \emph{precisión y recall} lo
son: si $P=0{,}9$ y $R=0{,}1$, $F1 = 2\cdot\frac{0{,}9\cdot0{,}1}{1{,}0}=0{,}18$ (no $0{,}5$).

\nivel{Anclaje en los datos}
En la Tabla~V, XLM-RoBERTa logra $F1_{\text{neu}}=0{,}500$ (el paper lo llama «equivalente a
asignación aleatoria»), BETO $0{,}550$, mientras pos/neg rondan $0{,}70$--$0{,}74$. LoRA cierra ese
hueco ($0{,}613$ el 4B, $0{,}631$ el 1.7B). Como macro-F1 da igual peso a la clase difícil, premia
ese cierre: por eso LoRA «gana» en macro-F1 aunque empate en las clases fáciles.

\begin{analisis}
  \item[Qué] Media no ponderada de la F1 de las tres clases; métrica primaria del estudio.
  \item[Por qué] El test está balanceado (290/clase) pero la \emph{dificultad} no; macro-F1 trata
        las clases por igual y refleja el rendimiento en la clase difícil.
  \item[Alternativas descartadas] \emph{Accuracy/micro-F1} (se dejan dominar por el agregado y
        ocultan el fallo en neutral); \emph{weighted-F1} (pondera por frecuencia; con clases
        balanceadas apenas cambia y no penaliza la clase difícil).
  \item[Qué significa] Macro-F1 alto = rendimiento \emph{equilibrado} en las tres clases.
  \item[Qué NO significa] No mide calibración; no distingue \emph{tipos} de error
        (pos$\leftrightarrow$neu vs pos$\leftrightarrow$neg); y \textbf{no es comparable entre
        datasets con distinto número de clases}.
\end{analisis}

\nivel{Autotest}
\begin{enumerate}[nosep]
  \item Define con tus palabras precisión y recall para la clase neutral.
  \item ¿Por qué F1 usa media armónica y no aritmética? Da un ejemplo numérico.
  \item Un modelo que predice \emph{siempre} «positivo» en test balanceado de 3 clases: ¿qué
        accuracy y qué macro-F1 aproximado tiene? (Pista: acc~$=1/3$; solo pos con recall~1 y
        precisión~$\approx1/3 \Rightarrow F1_{\text{pos}}\approx0{,}5$, las otras $0 \Rightarrow$
        macro~$\approx0{,}17$.)
  \item ¿Por qué el paper la llama «la cifra más informativa»? Relaciónalo con la Tabla~V.
\end{enumerate}

% ============================================================
\section{El objetivo SFT: solo la etiqueta cuenta}\label{sec:sft}
% ============================================================

\nivel{Definición rigurosa}
El \emph{fine-tuning} es supervisado (SFT) con \emph{cross-entropy}, pero \textbf{enmascarando el
prompt}: a todos los tokens del prompt se les asigna la etiqueta $-100$ (valor que la función de
pérdida ignora), de modo que solo el \textbf{token de etiqueta} (una de \emph{negative},
\emph{neutral}, \emph{positive}) y el token \texttt{EOS} contribuyen a la pérdida. Formalmente, si
$\mathcal{M}$ es el conjunto de posiciones no enmascaradas,
\[
\mathcal{L} = -\!\!\sum_{t \in \mathcal{M}} \log p_\theta\big(y_t \mid y_{<t}\big),
\qquad \mathcal{M} = \{\text{token de etiqueta},\ \texttt{EOS}\}.
\]
La estructura del prompt (system prompt, chat template, \emph{thinking mode} off) es \emph{idéntica}
a la del baseline de prompting.
\fuente{\S III-C «Model Selection and SFT Objective»; \texttt{configs/prompting\_protocol.yaml}}

\nivel{La intuición (apoyo)}
No queremos enseñar al modelo a \emph{recitar la pregunta}, solo a \emph{dar la respuesta correcta}.
Por eso «tapamos» (máscara $-100$) todo el enunciado y solo puntuamos la palabra final
(la etiqueta). Es como un examen donde solo corrigen la casilla de la respuesta, no si copiaste bien
el enunciado.

\nivel{Por qué importa este detalle}
Como el prompt es idéntico entre prompting y fine-tuning, y solo la etiqueta entra en la pérdida,
\textbf{la única variable que cambia entre «pedírselo» y «afinarlo» son los pesos actualizados}. Eso
es lo que hace \emph{limpio} el experimento de Cut~A: el cruce prompting$\to$fine-tuning no se puede
atribuir a diferencias de protocolo.

\begin{analisis}
  \item[Qué] Cross-entropy con enmascarado del prompt ($-100$); solo etiqueta + \texttt{EOS}
        contribuyen.
  \item[Por qué] Aislar la variable «pesos» entre prompting y fine-tuning, y centrar el aprendizaje
        en producir la etiqueta correcta.
  \item[Alternativas descartadas] Pérdida sobre \emph{todos} los tokens (mezclaría aprender la tarea
        con memorizar el prompt) o cabezal de clasificación dedicado (rompería la equivalencia con
        el harness generativo del baseline).
  \item[Qué significa] Prompting y fine-tuning comparten todo salvo los pesos; el experimento de
        cruce es justo.
  \item[Qué NO significa] No que el modelo aprenda «formato» (eso ya lo da el prompt); no que se
        entrene un clasificador clásico: sigue siendo generación, solo que supervisada en el token
        de etiqueta.
\end{analisis}

\nivel{Autotest}
\begin{enumerate}[nosep]
  \item ¿Qué representa el valor $-100$ y sobre qué tokens se aplica?
  \item ¿Por qué es importante que el prompt sea idéntico al del baseline de prompting?
  \item \textbf{Derivación en blanco.} Escribe la forma de $\mathcal{L}$ indicando explícitamente
        qué posiciones componen $\mathcal{M}$ y por qué el resto no suma a la pérdida.
\end{enumerate}

% ============================================================
\section{Caracterización en sigmas: magnitud, no significancia}\label{sec:sigma}
% ============================================================

\nivel{Definición rigurosa}
En Cut~C, cada celda (variedad origen $\to$ variedad destino) se evalúa con 3 semillas
$\{42,43,44\}$; se reporta media $\pm$ desviación típica (std) entre semillas. Un «gap» $\Delta$
entre dos celdas se expresa en \textbf{unidades de std de semilla}:
\[
\text{gap en } \sigma = \frac{|\Delta|}{\sigma_{\text{semilla}}}.
\]
Ejemplo (Qwen3-1.7B, CR$\to$ES): ES en-dominio $0{,}644\pm0{,}007$, CR$\to$ES $0{,}613\pm0{,}007$,
luego $\Delta=-0{,}031$, y $0{,}031/0{,}007\approx 4{,}4\sigma$.
\notrace{el paper reporta $4{,}6\sigma$ para esta celda; recalculando con la $\sigma$ redondeada de
la tabla sale $\approx4{,}4\sigma$ (diferencia por redondeo de $\sigma$ antes de dividir)}.
\fuente{\S IV-C de \texttt{paper/main.tex}; Tabla~III (\texttt{tab:cutC})}

\nivel{La intuición (apoyo)}
«$4{,}6\sigma$» suena a estadística, pero aquí significa solo: \emph{la caída es 4,6 veces el ruido
que hay entre semillas}. Es una forma de decir «esto es grande comparado con lo que varía al cambiar
la semilla aleatoria», no «esto es significativo» en el sentido de un test de hipótesis.

\nivel{El resultado y su honestidad}
La penalización al evaluar en ES (entrenando en CR/PE) es $0{,}031$--$0{,}039$ de F1
($2{,}6$--$4{,}6\sigma$), y \textbf{replica en 1.7B y 4B}. El paper avisa explícitamente: no hay test
de hipótesis (permutación, $t$ de Student), ni corrección por comparaciones múltiples sobre la
matriz $3\times3\times2$; el criterio de $2\sigma$ es un \emph{heurístico de magnitud}, no un
$p$-valor.

\begin{analisis}
  \item[Qué] Expresar los gaps de transferencia en múltiplos de la std entre semillas.
  \item[Por qué] Con corpus pequeños (539--738 ejemplos) y 3 semillas, da una medida \emph{honesta}
        de magnitud relativa al ruido, sin sobrevender significancia estadística.
  \item[Alternativas descartadas] Test formal con corrección de comparaciones múltiples (queda como
        \emph{future work}; con 3 semillas tendría poca potencia); reportar solo $\Delta$ sin
        contexto de ruido (no diría si es grande o pequeño).
  \item[Qué significa] La asimetría hacia ES es \emph{grande respecto al ruido} y \emph{reproducible}
        en ambos tamaños de modelo.
  \item[Qué NO significa] No es significancia estadística ($p$-valor); no prueba causalidad
        dialectal; con más semillas/datos podría matizarse. La excepción CR$\to$PE (3,9$\sigma$ en
        1.7B) \emph{no} replica en 4B (0,3$\sigma$): solo la columna ES es robusta.
\end{analisis}

\nivel{Autotest}
\begin{enumerate}[nosep]
  \item ¿Qué afirma y qué \emph{no} afirma un «gap de $4{,}6\sigma$» aquí?
  \item ¿Por qué el paper no aplica un test de hipótesis formal? ¿Qué tendría que cambiar para
        hacerlo?
  \item \textbf{Derivación en blanco.} Dados ES en-dominio $0{,}644\pm0{,}007$ y PE$\to$ES
        $0{,}613\pm0{,}008$, calcula $\Delta$ y el gap en $\sigma$. (Solución: $\Delta=-0{,}031$;
        $0{,}031/0{,}008\approx 3{,}9\sigma$, coherente con el $3{,}8\sigma$ del paper.)
\end{enumerate}

% ============================================================
\section{El punto de cruce y su mecanismo}\label{sec:crossover}
% ============================================================

\nivel{Definición rigurosa}
El \textbf{punto de cruce} es el menor $n$ (tamaño del set de entrenamiento) para el que LoRA supera
el \emph{mejor} resultado de prompting del \emph{mismo} modelo (N-independiente). De la Tabla~II:
\begin{itemize}[nosep]
  \item Qwen3-1.7B: mejor prompting $0{,}560$ ($k{=}8$). LoRA está por debajo en $n{=}4$
        ($0{,}539\pm0{,}010$) y \textbf{lo supera en $n{=}7$} ($0{,}590\pm0{,}063$).
  \item Qwen3-4B: 0-shot ya $0{,}633$; mejor prompting $0{,}652$ ($k{=}4$). En $n{=}7$ ($0{,}645\pm0{,}038$)
        y $n{=}25$ ($0{,}645\pm0{,}037$) sigue por debajo; \textbf{cruza hacia $n\approx50$}
        ($0{,}666\pm0{,}034$).
\end{itemize}
\fuente{\S IV-A y Tabla~II (\texttt{tab:cutA}); mecanismo en \S II-C y \S V}

\nivel{La intuición (apoyo) y el mecanismo}
«Pedírselo» (prompting) tiene un \textbf{suelo de calidad} que depende de lo que el modelo ya sabe
de fábrica. Las demostraciones del prompt sobre todo comunican \emph{formato} y \emph{espacio de
etiquetas}, no enseñan el mapeo correcto \externo{Min et al., 2022}. Por eso un modelo más grande
(4B) parte de un suelo más alto ($0{,}652$ vs $0{,}560$) y el fine-tuning necesita \emph{más señal}
(más ejemplos) para superarlo. De ahí que el cruce se desplace a la derecha con el tamaño: no hay un
umbral universal.

\begin{analisis}
  \item[Qué] El cruce prompting$\to$fine-tuning depende del tamaño: $<7$ ejemplos (1.7B) vs
        $\approx50$ (4B).
  \item[Por qué] Suelo de prompting $\propto$ conocimiento preentrenado; suelo alto $\Rightarrow$
        más datos para superarlo.
  \item[Alternativas descartadas] Buscar un umbral único universal (no existe; varía con el suelo
        del modelo, y probablemente con tarea/dominio).
  \item[Qué significa] Regla práctica: con muy pocos labels ($n<10$) y modelo pequeño, afina ya; con
        4B, comprueba primero si el suelo de prompting te basta, porque afinar solo gana fiable hacia
        $n\approx50$.
  \item[Qué NO significa] Los números ($n{<}7$, $n{\approx}50$) \emph{no} son constantes
        universales; lo reproducible es la \emph{existencia} del desplazamiento con el tamaño.
\end{analisis}

\nivel{Autotest}
\begin{enumerate}[nosep]
  \item ¿Por qué un modelo más grande \emph{retrasa} el punto de cruce?
  \item ¿Qué papel juegan las demostraciones del prompt según Min et al.?
  \item ¿Qué parte de este hallazgo es universal y qué parte depende del modelo?
\end{enumerate}

% ============================================================
\section{Muestreo estratificado y disjunto por semilla}\label{sec:subsampling}
% ============================================================

\nivel{Definición rigurosa}
Para Cut~A se construyen subconjuntos de tamaño $n\in\{4,7,10,16,25,50,100,250,500,1000\}$ más el
total ($1839$). Cada semilla $\{42,43,44\}$ extrae un subconjunto \textbf{estratificado}
(las 3 clases presentes en cada celda; en $n{=}4$: neg$=2$, neu$=1$, pos$=1$) e
\textbf{independiente/disjunto} de los de las otras semillas (0 ejemplos compartidos en $n{=}4$, 0 en
$n{=}10$).
\fuente{\S III-A; \S IV-A; \texttt{scripts/build\_fraction\_subsets.py}}

\nivel{La intuición (apoyo)}
Queremos que la barra de error (la std) signifique algo real. Si las tres semillas vieran \emph{los
mismos} ejemplos y solo cambiara la inicialización, la std mediría «ruido de pesos». Al hacer que
cada semilla vea ejemplos \emph{distintos}, la std mide algo más informativo: \textbf{cuánto depende
el resultado de \emph{qué} datos te tocaron} (varianza de muestreo de datos). La estratificación
evita que, por azar, una celda diminuta se quede sin alguna clase.

\begin{analisis}
  \item[Qué] Submuestras estratificadas e independientes por semilla.
  \item[Por qué] Que la std refleje varianza de \emph{muestreo de datos}, no solo de inicialización;
        garantizar las 3 clases en celdas diminutas.
  \item[Alternativas descartadas] Submuestras solapadas / fijas entre semillas (la std subestimaría
        la incertidumbre real); muestreo no estratificado (riesgo de clases ausentes en $n$ pequeños).
  \item[Qué significa] Las barras de error de Cut~A son honestas respecto a la sensibilidad a los
        datos.
  \item[Qué NO significa] No elimina toda incertidumbre: con $n{=}4$ la varianza es enorme
        ($\pm0{,}089$ en 4B) precisamente porque depende muchísimo de qué 4 ejemplos tocaron.
\end{analisis}

\nivel{Autotest}
\begin{enumerate}[nosep]
  \item ¿Qué mide la std en Cut~A gracias a que las submuestras son disjuntas?
  \item ¿Por qué se estratifica en lugar de muestrear al azar puro?
  \item ¿Por qué la varianza es máxima en $n{=}4$ y se contrae al crecer $n$?
\end{enumerate}

% ============================================================
\section{Frontera de Pareto y dominación}\label{sec:pareto}
% ============================================================

\nivel{Definición rigurosa}
Un método \textbf{domina} a otro si es mejor o igual en \emph{todos} los ejes y estrictamente mejor
en al menos uno. La \textbf{frontera de Pareto} es el conjunto de métodos no dominados. En Cut~B los
ejes son calidad (macro-F1) y coste (parámetros entrenables, VRAM, tiempo, latencia). LoRA domina a
los encoders: más F1 ($0{,}707/0{,}698$ vs $0{,}661/0{,}646$) con $6$--$12$~M parámetros entrenables.
\fuente{\S IV-B; Tabla~IV (\texttt{tab:cutB}); Fig.~2}

\nivel{La intuición (apoyo)}
Imagina comprar coche por precio y consumo: un coche está «en la frontera» si no hay otro más barato
\emph{y} que gaste menos a la vez. Los que sí tienen a alguien mejor en todo quedan «dominados».

\begin{analisis}
  \item[Qué] LoRA ocupa la frontera calidad-coste; full-FT y encoders quedan dominados (salvo full-FT
        en latencia).
  \item[Por qué] Mejor calidad a una fracción del coste de parámetros/VRAM.
  \item[Alternativas descartadas] Mirar solo calidad (ocultaría el coste) o solo coste (ocultaría la
        calidad): la gracia es el \emph{trade-off} conjunto.
  \item[Qué significa] Para presupuesto limitado, no hay escenario donde full-FT de un encoder gane a
        LoRA en los ejes medidos.
  \item[Qué NO significa] \textbf{No} que LoRA tenga la mayor F1 absoluta del paper: RoBERTuito saca
        $0{,}758$, pero está \emph{fuera} de la frontera por ser out-of-domain (TASS~2020, zero-shot)
        y no comparable de forma controlada (esto es la debilidad W1).
\end{analisis}

\nivel{Autotest}
\begin{enumerate}[nosep]
  \item Define «dominar» con tus palabras y di por qué full-FT queda dominado salvo en latencia.
  \item ¿Por qué RoBERTuito (F1 $0{,}758$) no está en la frontera pese a tener la F1 más alta?
\end{enumerate}

% ============================================================
\section{Los cuatro ejes de coste}\label{sec:costaxes}
% ============================================================

\nivel{Definición rigurosa}
Por cada celda experimental se miden cuatro costes: (1)~\textbf{parámetros entrenables} (los que se
actualizan, no los totales); (2)~\textbf{VRAM pico} en GB, vía
\texttt{torch.cuda.max\_memory\_allocated()}; (3)~\textbf{tiempo de entrenamiento} (wall-clock, s);
(4)~\textbf{latencia} de inferencia (ms/ejemplo, promediada sobre el test).
\fuente{\S III-D; Tabla~IV}

\nivel{La intuición (apoyo)}
«Calidad» a secas engaña: un modelo carísimo que acierta un poco más no siempre compensa. Estos
cuatro ejes responden a preguntas reales de despliegue: ¿qué guardo y reentreno por tarea? (params)
¿me cabe en la GPU? (VRAM) ¿cuánto tardo en entrenar? (tiempo) ¿responde rápido? (latencia).

\begin{analisis}
  \item[Qué] Cuatro dimensiones de coste medidas por celda junto a la calidad.
  \item[Por qué] Hacer comparables modelos generativos y encoders en \emph{coste}, no solo en F1;
        validar viabilidad en T4 16~GB (la VRAM pico es la cota clave).
  \item[Alternativas descartadas] Reportar calidad aislada del coste (lo que hacía la literatura
        previa, según \S II); usar parámetros \emph{totales} en vez de entrenables (no refleja el
        coste real de PEFT por tarea).
  \item[Qué significa] La distinción entrenables vs totales es la que justifica el ahorro de PEFT
        (6--12~M frente a backbones de miles de millones).
  \item[Qué NO significa] La VRAM se midió en H200, pero es una \emph{cota superior} válida para T4;
        los \emph{tiempos} sí cambiarían en T4 (ver debilidad W4).
\end{analisis}

\nivel{Autotest}
\begin{enumerate}[nosep]
  \item ¿Por qué se miden parámetros \emph{entrenables} y no totales?
  \item ¿Cuál de los cuatro ejes es el realmente relevante para la afirmación «free-tier / T4» y por
        qué?
\end{enumerate}

% ============================================================
\section{El suelo de pasos y el calendario de entrenamiento}\label{sec:maxsteps}
% ============================================================

\nivel{Definición rigurosa}
Para que las fracciones diminutas no queden infraentrenadas, el número de pasos de optimización se
acota por abajo:
\[
\text{max\_steps} = \max\!\left(80,\ 5\cdot\Big\lceil \tfrac{n}{8} \Big\rceil\right).
\]
Calendario común a LoRA/QLoRA: lr~$=2\times10^{-4}$, scheduler \emph{cosine} con warmup ratio
$0{,}1$, 5 épocas, batch 8, grad.\ accum.\ 1, longitud máxima 256, bf16. La val macro-F1 se evalúa 4
veces durante el entrenamiento.
\fuente{\S III-C «Training Schedule»; \texttt{configs/sprint3/grid/lora\_qwen1.7b\_n100\_s42.yaml}}

\nivel{La intuición (apoyo)}
Con muy pocos ejemplos, «5 épocas» son poquísimos pasos: el modelo apenas se entera. El suelo de
80 pasos garantiza un mínimo de práctica aunque solo haya 4 ejemplos. Por encima de cierto $n$, manda
la fórmula $5\lceil n/8\rceil$ (las 5 épocas reales).

\begin{analisis}
  \item[Qué] Suelo de 80 pasos; resto del calendario fijo.
  \item[Por qué] Evitar que $n{=}4,7,10$ queden infraentrenados y contaminen la curva de Cut~A.
  \item[Alternativas descartadas] Fijar «5 épocas» sin suelo (las celdas diminutas casi no se
        entrenarían) o un nº de pasos fijo igual para todo $n$ (sobreentrenaría los $n$ grandes).
  \item[Qué significa] La curva F1-vs-$n$ es comparable entre tamaños porque ninguna celda está
        infraentrenada por accidente.
  \item[Qué NO significa] No que 80 pasos basten para \emph{aprender} con 4 ejemplos: la enorme
        varianza en $n{=}4$ muestra que el límite es el dato, no los pasos.
\end{analisis}

\nivel{Autotest}
\begin{enumerate}[nosep]
  \item ¿Qué problema evita el suelo de 80 pasos?
  \item \textbf{Derivación en blanco.} Calcula max\_steps para $n{=}4$ y para $n{=}1000$.
        (Solución: $n{=}4$: $\max(80,5\lceil4/8\rceil)=\max(80,5)=80$; $n{=}1000$:
        $\max(80,5\lceil1000/8\rceil)=\max(80,5\cdot125)=625$.)
\end{enumerate}

% ============================================================
\section{Selección del mejor checkpoint (val macro-F1)}\label{sec:checkpoint}
% ============================================================

\nivel{Definición rigurosa}
No se usa el \emph{último} paso: se selecciona el checkpoint con la \textbf{mayor macro-F1 de
validación} (implementado en \texttt{BestValF1Callback}). La validación se evalúa 4 veces durante el
entrenamiento; la selección de hiperparámetros usa exclusivamente el split de validación (324
ejemplos); el test (870) se toca una sola vez, al final.
\fuente{\S III-C «Model Selection»; \S III-D}

\nivel{La intuición (apoyo)}
Entrenar de más puede empeorar (overfitting). En vez de quedarte con «la última foto», te quedas con
«la mejor foto» según un examen de práctica (validación) que no es el examen final (test). Así no
contaminas el test ni te quedas con un modelo pasado de cocción.

\begin{analisis}
  \item[Qué] Elegir el checkpoint de mejor val-F1, no el final.
  \item[Por qué] Evitar overfitting tardío y respetar la separación validación/test (sin fuga).
  \item[Alternativas descartadas] Usar el último paso (riesgo de overfitting); seleccionar sobre el
        test (fuga de datos, invalidaría la comparación).
  \item[Qué significa] Las cifras de test son honestas: el test se usó una sola vez.
  \item[Qué NO significa] No garantiza el óptimo global; solo el mejor de las 4 evaluaciones de
        validación realizadas.
\end{analisis}

\nivel{Autotest}
\begin{enumerate}[nosep]
  \item ¿Por qué no quedarse con el último checkpoint?
  \item ¿Por qué seleccionar sobre validación y no sobre test?
\end{enumerate}

% ============================================================
\section{Desbalance de clases y el descarte de NONE (Cut C)}\label{sec:none}
% ============================================================

\nivel{Definición rigurosa}
InterTASS~2018 anota 4 polaridades (P, N, NEU, NONE). Se mapean P$\to$positivo, N$\to$negativo,
NEU$\to$neutral, y se \textbf{descarta NONE} («sin opinión / sin objetivo»). Tras el descarte: 738~ES
/ 539~CR / 541~PE de entrenamiento; hold-out estratificado del 15\% (131/96/96) para selección;
evaluación en los dev splits (444/242/262). La clase neutral es más frecuente en PE ($\approx26\%$)
que en ES/CR ($\approx15\%$).
\fuente{\S III-A; \S IV-C}

\nivel{La intuición (apoyo)}
NONE no es «neutral»: neutral es «tiene postura, pero ni positiva ni negativa»; NONE es «no hay
objetivo del que opinar». Mezclarlos sería meter peras en el cajón de las manzanas y \emph{contaminar}
el esquema de polaridad. Y como PE tiene casi el doble de neutrales, su distribución es más
desbalanceada, lo que explica más varianza y peor F1 en neutral allí.

\begin{analisis}
  \item[Qué] Descartar NONE; trabajar con 3 clases; convivir con desbalance (PE 26\% neu).
  \item[Por qué] NONE es conceptualmente distinto de neutral; fusionarlo contaminaría la polaridad.
  \item[Alternativas descartadas] Mapear NONE$\to$neutral (contaminación conceptual); descartar
        también NEU (perdería una de las 3 clases del estudio).
  \item[Qué significa] El esquema de 3 clases de Cut~C es coherente con el de Cuts~A--B.
  \item[Qué NO significa] No elimina el desbalance: PE sigue teniendo neutral sobre-representado, con
        efectos en su F1 y varianza (limitación declarada).
\end{analisis}

\nivel{Autotest}
\begin{enumerate}[nosep]
  \item ¿Por qué NONE no se mapea a neutral?
  \item ¿Qué consecuencia tiene que PE tenga 26\% de neutral frente al 15\% de ES/CR?
\end{enumerate}

% ============================================================
\section{Por qué la clase neutral es la difícil}\label{sec:neutral}
% ============================================================

\nivel{Definición rigurosa}
En la Tabla~V, el \emph{ranking} por clase es idéntico en \emph{todos} los modelos y métodos:
positivo $\geq$ negativo $>$ neutral. Neutral es siempre la más débil (encoders $0{,}500$--$0{,}550$;
LoRA $0{,}613$--$0{,}631$). Que el orden se repita en encoders y generativos indica que la dificultad
de neutral es \textbf{propiedad de la tarea}, no artefacto de una arquitectura.
\fuente{\S IV-B «Per-Class Error Analysis»; Tabla~V}

\nivel{La intuición (apoyo)}
Positivo y negativo tienen \emph{señal} (palabras de alegría o enfado). Neutral es la \emph{ausencia}
de señal afectiva: hay que detectar que \emph{no} hay emoción, lo cual no tiene un patrón superficial
distintivo (puede ser lenguaje coloquial sin carga, hechos, u opiniones que se cancelan). Detectar un
«no-algo» es intrínsecamente más difícil que detectar un «algo».

\begin{analisis}
  \item[Qué] Neutral es la clase más difícil para todos; LoRA la mejora más que los encoders.
  \item[Por qué] Neutral = ausencia de señal afectiva, sin patrón superficial.
  \item[Alternativas descartadas] Atribuirlo a la arquitectura encoder (descartado: el ranking se
        repite también en generativos).
  \item[Qué significa] La ventaja de LoRA en neutral sugiere que el preentrenamiento más amplio de
        Qwen aporta señal extra para identificar el «no-sentimiento».
  \item[Qué NO significa] Esa explicación causal (Qwen «sabe más» que BETO) es una \emph{hipótesis
        razonada}, no demostrada (debilidad W9); el dato duro es el $+0{,}06$/$+0{,}11$ en neutral,
        no su causa.
\end{analisis}

\nivel{Autotest}
\begin{enumerate}[nosep]
  \item ¿Por qué detectar «neutral» es más difícil que detectar positivo o negativo?
  \item ¿Qué evidencia indica que es propiedad de la tarea y no de la arquitectura?
\end{enumerate}

% ============================================================
\section{La familia PEFT: de dónde viene LoRA y qué hay más allá}\label{sec:peftfamily}
% ============================================================

\nivel{Definición rigurosa}
PEFT = \emph{Parameter-Efficient Fine-Tuning}: confinar la adaptación a un subconjunto pequeño de
parámetros con el backbone congelado. Hitos: \textbf{adapters} \externo{Houlsby et al., 2019}
(capas cuello de botella entre bloques; $<4\%$ de parámetros, pero añaden latencia);
\textbf{LoRA} \externo{Hu et al., 2021} (sin latencia al fundir); \textbf{QLoRA}
\externo{Dettmers et al., 2023} (LoRA + cuantización NF4); \textbf{DoRA} \externo{Liu et al., 2022}
(descompone la actualización en magnitud y dirección). Todo unificado en la librería
\textbf{PEFT} \externo{Mangrulkar et al., 2022}.
\fuente{\S II-A; \S VI (Future Work) de \texttt{paper/main.tex}}

\nivel{La intuición (apoyo)}
Hay muchas formas de «pegar notas adhesivas» a un modelo congelado. Los adapters clásicos añaden
piezas nuevas (y por eso pesan al inferir). LoRA logra el mismo efecto sin coste de inferencia
(se funde). DoRA e IA$^3$ son primos más recientes que podrían exprimir aún más el trade-off
calidad-por-parámetro: por eso aparecen como \emph{future work}, no como parte de los experimentos.

\begin{analisis}
  \item[Qué] Situar LoRA/QLoRA dentro de la familia PEFT y sus alternativas.
  \item[Por qué] Justificar la elección (LoRA: sin latencia al fundir; QLoRA: memoria) y delimitar el
        alcance (DoRA/IA$^3$ no evaluados).
  \item[Alternativas descartadas] Adapters clásicos (latencia en inferencia); DoRA/IA$^3$ (fuera del
        alcance, \emph{future work}).
  \item[Qué significa] LoRA/QLoRA son una elección informada dentro de un abanico, no la única
        opción.
  \item[Qué NO significa] No que LoRA sea óptimo entre todos los PEFT; DoRA/IA$^3$ podrían mejorar la
        frontera (no testado aquí).
\end{analisis}

\nivel{Autotest}
\begin{enumerate}[nosep]
  \item ¿Qué diferencia a los adapters clásicos de LoRA en \emph{inferencia}?
  \item ¿Por qué DoRA e IA$^3$ aparecen como \emph{future work} y no en los experimentos?
\end{enumerate}

\end{document}
